\chapter{Estado del arte}\label{cap:estado}

El análisis computacional de la música ofrece un marco cuantificable para estudiar la melodía, la armonía, el ritmo y el
timbre, y resulta especialmente útil frente a músicas que desafían los modelos teóricos convencionales. La
microtonalidad, el uso de intervalos menores que el semitono, es uno de esos casos: la notación occidental y la mayoría
de las herramientas de análisis no representan con precisión las alturas que caen entre los semitonos.

Este trabajo toma elementos de seis líneas de investigación. La tabla~\ref{tab:estado} resume, para cada una, qué
resuelve, qué le falta para esta música y qué se usa aquí.

\begin{table}[ht]
  \centering\small
  \caption{Líneas de trabajo previo y su relación con este trabajo.}\label{tab:estado}
  \begin{tabularx}{\textwidth}{@{}>{\raggedright}p{2.2cm}>{\raggedright}p{3.1cm}>{\raggedright}X>{\raggedright}X>{\raggedright\arraybackslash}X@{}}
    \toprule
    Línea & Trabajos & Qué resuelven & Qué les falta aquí & En este trabajo \\
    \midrule
    Altura desde el audio & \citet{mauch2014}; \citet{bittner2022} & Melodía y notas a partir de la grabación &
      pYIN sigue una sola voz; Basic Pitch da notas MIDI enteras y el cuarto de tono queda como desafinación &
      Saliencia armónica y consenso entre vueltas del looper \\
    Separación de fuentes & \citet{rouard2023}; \citet{lu2024} & Aislar bajo, guitarra y batería &
      No transcriben: solo preparan la señal & Demucs para guitarra y bajo \\
    MIR microtonal & \citet{bozkurt2008}; \citet{gedik2010} & Histogramas de altura finos para el makam turco &
      Música de una sola voz, sin acordes ni looper & Rejilla de 24-EDO con casillas de $\pm25$ cents \\
    Armonía funcional & \citet{riemann1893}; \citet{lendvai1971}; \citet{piston1987} &
      Función de cada acorde: tónica, subdominante o dominante & Pensada en 12 notas; Piston deja fuera lo microtonal &
      Calidades, funciones y ejes extendidos a 24-EDO \\
    Estilo con corpus & \citet{declercq2011}; \citet{broze2013} & Perfiles de armonía del rock y del jazz &
      Datos simbólicos en 12 notas & Perfiles de 24 clases, Jensen-Shannon y percentil \\
    Métrica irregular & \citet{holzapfel2014}; \citet{srinivasamurthy2016} & Compases no occidentales desde el audio &
      Sin la dimensión de altura & Pulsos y compases de 2 a 8 pulsos \\
    \bottomrule
  \end{tabularx}
\end{table}

\section{Altura desde el audio}

Los estimadores de frecuencia fundamental como YIN \citep{decheveigne2002} y pYIN \citep{mauch2014} siguen con precisión
una línea monofónica y entregan alturas continuas que pueden expresarse como MIDI fraccionario o desviación en cents. Los
modelos de redes neuronales, como CREPE \citep{kim2018}, mejoran la estimación en una sola voz. Para música polifónica, la
extracción de melodía basada en contornos de saliencia \citep{salamon2012} y los transcriptores multinota, como Basic Pitch
\citep{bittner2022}, entregan varias notas a la vez, pero cuantizadas a doce notas: un cuarto de tono les sale como una
nota desafinada.

\section{Separación de fuentes}

Demucs \citep{rouard2023} y BS-RoFormer \citep{lu2024} separan una mezcla en pistas de instrumentos (bajo, batería, voz,
guitarra y otros). Son un paso previo útil, pero no transcriben. REPET \citep{rafii2013} separa lo que se repite de lo
que no, una idea cercana al funcionamiento de un looper.

\section{MIR microtonal}

El estudio computacional del makam turco ha desarrollado histogramas de altura de alta resolución y comparaciones entre
modelos teóricos de afinación y mediciones \citep{bozkurt2008,bozkurt2009,gedik2010,bozkurt2014}. Estas técnicas trabajan
con música de una sola voz (heterofonía, en el mejor de los casos), sin acordes y sin la estructura en capas de un
looper. \citet{serra2011} propuso un enfoque multicultural para la investigación en información musical: los modelos
del área se construyeron sobre la música occidental y tienden a llevar ese sesgo a otras tradiciones.

\section{Armonía funcional}

Las teorías de \citet{riemann1893}, \citet{lendvai1971} y \citet{piston1987} clasifican los acordes por su función
tonal. Todas están pensadas para doce notas, y Piston excluye explícitamente las escalas microtonales.

\section{Estilo con corpus}

Los estudios de corpus describen estadísticamente la armonía de un estilo: el rock \citep{declercq2011,temperley2013},
la práctica común \citep{temperley2009} y el jazz \citep{broze2013,shanahan2012}. La comparación de distribuciones con
medidas de teoría de la información tiene base en la cognición musical \citep{krumhansl1990,huron2006,pearce2018}. Todos
estos corpus son simbólicos y están en doce notas.

\section{Métrica irregular}

El seguimiento de compases no occidentales desde el audio, como los usul turcos o los tala indios, se ha estudiado con
modelos bayesianos \citep{holzapfel2014,srinivasamurthy2016}, y el seguimiento de pulsos por programación dinámica es un
método clásico \citep{ellis2007}. Estos trabajos no incorporan la dimensión de altura.

\section{El vacío}

Ningún trabajo consultado reúne a la vez audio polifónico, 24-EDO, armonía funcional y comparación con corpus. Este
trabajo une las seis líneas para estudiar una banda de rock con cuartos de tono y looper.
